\section{Weighted Windows at Split Finite Places}
\label{fw:section}

We now place locally constant weights at the selected finite places.
This refines the choice of ideals in Lemma~\ref{geo:selection} in the
same way that the archimedean profiles refine a ball. We keep the
hypotheses (G1)--(G3) and the notation of Section~\ref{geo:section},
and put
\[
 \covol_0=\covol\bigl((\mathcal O_K)_\infty\bigr)=(\lambda/2)^d .
\]

\subsection{The Local Functional}

Let $L$ be a nonarchimedean local field with valuation ring $\mathcal
O$, uniformizer $\varpi$ and residue field of cardinality $Q$. We use
the additive Haar measure on $L$ with $|\mathcal O|=1$, the product
measure on $L^2$, and the multiplicative Haar measure $d\omega$ on
$\mathcal O^\times$ of total mass one. For $n\in\Z$ and
$\omega\in\mathcal O^\times$ put
\[
 s(n,\omega)=(\varpi^n\omega,\ \varpi^{-n}\omega^{-1})\in L^2,\qquad
 Tg(z)=\sum_{n\in\Z}\int_{\mathcal O^\times}g(z+s(n,\omega))\,d\omega .
\]
The two coordinates of $s(n,\omega)$ have product one; they model the
finite coordinates of an element of relative norm one at a split place.
For a nonnegative, locally constant, compactly supported function $g$
on $L^2$ with $g\ne0$, and $0<\delta<1$, $p=2/(1+\delta)$, put
\begin{equation}\label{fw:functional}
 A(g)=\int_{L^2}g^p,\qquad Z(g)=\int_{L^2}g\,Tg,\qquad
 \mathcal F_\delta(g)=\frac{Z(g)}{A(g)^{1+\delta}} .
\end{equation}
Only finitely many $n$ contribute to $Z(g)$: if $g(z)g(z+s(n,\omega))\ne0$
then both coordinates of $s(n,\omega)$ lie in the compact set
$\operatorname{supp}g-\operatorname{supp}g$, which bounds $n$ from above
and from below.

For $i\in\Z$ let $\mathcal B_i=\varpi^{-i}\mathcal O$, a ball of measure
$Q^i$. Write $[x]_+=\max(x,0)$.

\begin{lemma}\label{fw:gram-lemma}
For all integers $i_1,j_1,i_2,j_2$,
\begin{equation}\label{fw:rectangle-gram}
 \int_{L^2}\mathbf 1_{\mathcal B_{i_1}\times\mathcal B_{j_1}}\,
 T\mathbf 1_{\mathcal B_{i_2}\times\mathcal B_{j_2}}
 =Q^{\min(i_1,i_2)+\min(j_1,j_2)}\,
 [\max(i_1,i_2)+\max(j_1,j_2)+1]_+.
\end{equation}
\end{lemma}

\begin{proof}
Two balls of an ultrametric space are disjoint or nested. Hence
$\mathcal B_{i_1}\cap(\mathcal B_{i_2}-x)$ has measure
$Q^{\min(i_1,i_2)}$ if $x\in\mathcal B_{\max(i_1,i_2)}$, and is empty
otherwise. The left side of \eqref{fw:rectangle-gram} is therefore
$Q^{\min(i_1,i_2)+\min(j_1,j_2)}$ times the measure of the set of
$(n,\omega)$ with $\varpi^n\omega\in\mathcal B_{\max(i_1,i_2)}$ and
$\varpi^{-n}\omega^{-1}\in\mathcal B_{\max(j_1,j_2)}$. These conditions
say $-\max(i_1,i_2)\le n\le\max(j_1,j_2)$ and do not involve $\omega$,
and the unit integral has mass one.
\end{proof}

Put $\mathcal S_0=\mathcal O$ and $\mathcal S_i=\mathcal B_i\setminus
\mathcal B_{i-1}$ for $i\ge1$, the \emph{shells}, of measures $m_0=1$
and $m_i=Q^i-Q^{i-1}$.

\begin{definition}\label{fw:shell-profile}
Let $k\ge0$ be an integer, and let $w=(w_{ij})_{i,j\ge0}$ be nonnegative
weights, finitely many of them nonzero, with $w_{00}>0$. The
\emph{shell profile} of type $(Q,k,w)$ on $L^2$ is
\[
 g=\sum_{i,j\ge0}w_{ij}\mathbf 1_{\mathcal S_i\times\varpi^{-k}\mathcal S_j}.
\]
\end{definition}

A shell profile is nonnegative, locally constant and compactly
supported. It is invariant under $(x,y)\mapsto(\omega x,\omega'y)$ for
all $\omega,\omega'\in\mathcal O^\times$, since multiplication by a unit
preserves valuations and hence each shell; in particular it is even. It
is also invariant under translation by the group
\begin{equation}\label{fw:hard-period}
 \mathcal C=\mathcal O\times\varpi^{-k}\mathcal O,\qquad |\mathcal C|=Q^k,
\end{equation}
whatever the number of shells: adding an element of $\mathcal O$ does
not change the shell of a point. We call $\mathcal C$ the \emph{period}
of $g$. Define the symmetric matrix $\Delta=(\Delta_{ii'})_{i,i'\ge0}$ by
\[
 \Delta_{ii'}=m_{\min(i,i')}\ (i\ne i'),\qquad \Delta_{00}=0,\qquad
 \Delta_{ii}=im_i-Q^{i-1}\ (i\ge1),
\]
and the kernel
\[
 \mathcal K_{(i,j),(i',j')}=(k+1)m_im_j\mathbf 1_{i=i',\,j=j'}
 +m_j\Delta_{ii'}\mathbf 1_{j=j'}+m_i\Delta_{jj'}\mathbf 1_{i=i'}.
\]

\begin{lemma}\label{fw:shell-lemma}
The shell profile $g$ of type $(Q,k,w)$ satisfies
\begin{equation}\label{fw:shell-kernel}
 A(g)=Q^k\sum_{i,j}m_im_jw_{ij}^p,\qquad
 Z(g)=Q^k\sum_{(i,j),(i',j')}w_{ij}\,\mathcal K_{(i,j),(i',j')}\,w_{i'j'} .
\end{equation}
In particular $\mathcal F_\delta(g)$ depends only on the type, and the
hard window $w_{00}=1$, $w_{ij}=0$ otherwise, has
$\mathcal F_\delta(g)=(k+1)Q^{-\delta k}$.
\end{lemma}

\begin{proof}
The formula for $A(g)$ holds because the sets
$\mathcal S_i\times\varpi^{-k}\mathcal S_j$ are disjoint of measure
$Q^km_im_j$. For $Z(g)$, write $\mathbf 1_{\mathcal S_0}=\mathbf
1_{\mathcal B_0}$ and $\mathbf 1_{\mathcal S_i}=\mathbf 1_{\mathcal
B_i}-\mathbf 1_{\mathcal B_{i-1}}$ for $i\ge1$, and note
$\varpi^{-k}\mathcal B_j=\mathcal B_{j+k}$. Only balls with nonnegative
indices occur, so by \eqref{fw:rectangle-gram} and $k\ge0$ the positive
part can be dropped, and the Gram value of $\mathcal B_{i_1}\times
\mathcal B_{j_1+k}$ and $\mathcal B_{i_2}\times\mathcal B_{j_2+k}$, for
$i_1,j_1,i_2,j_2\ge0$, is
\[
 Q^k\bigl[(k+1)Q^{\min(i_1,i_2)}Q^{\min(j_1,j_2)}
 +\max(i_1,i_2)Q^{\min(i_1,i_2)}Q^{\min(j_1,j_2)}
 +Q^{\min(i_1,i_2)}\max(j_1,j_2)Q^{\min(j_1,j_2)}\bigr].
\]
For a function $\varphi$ of two nonnegative integers put
$(\partial\varphi)(i,i')=\varphi(i,i')-\varphi(i-1,i')-\varphi(i,i'-1)
+\varphi(i-1,i'-1)$, where terms with an index $-1$ are omitted; this is
the expansion of the shells into balls. Apply $\partial$ in
$(i_1,i_2)$ and in $(j_1,j_2)$. A direct computation gives, for
$i,i'\ge0$, that $\partial Q^{\min(i_1,i_2)}$ at $(i,i')$ is
$m_i\mathbf 1_{i=i'}$, and that
$\partial\bigl(\max(i_1,i_2)Q^{\min(i_1,i_2)}\bigr)$ is $\Delta_{ii'}$.
For example, when $1\le i<i'$ the latter is
$i'Q^i-i'Q^{i-1}-(i'-1)Q^i+(i'-1)Q^{i-1}=m_i$, and when $i=i'\ge1$ it is
$iQ^i-2iQ^{i-1}+(i-1)Q^{i-1}=im_i-Q^{i-1}$. This gives the formula for
$Z(g)$. For the hard window only $\mathcal K_{(0,0),(0,0)}=k+1$
contributes.
\end{proof}

All entries of $\mathcal K$ are nonnegative, since
$\Delta_{ii}=Q^{i-1}(iQ-i-1)\ge0$ for $i\ge1$. Hence
$Z(g)\ge(k+1)Q^kw_{00}^2>0$ for every shell profile.

For product weights $w_{ij}=x_ix'_j$, with vectors $x,x'$ of
nonnegative numbers, put
\[
 P_x=\sum_im_ix_i^p,\qquad S_x=\sum_im_ix_i^2,\qquad
 R_x=\sum_{i,i'}x_i\Delta_{ii'}x_{i'},
\]
and define $P_{x'},S_{x'},R_{x'}$ in the same way. Then
\eqref{fw:shell-kernel} becomes
\begin{equation}\label{fw:product-shells}
 A(g)=Q^kP_xP_{x'},\qquad
 Z(g)=Q^k\bigl[(k+1)S_xS_{x'}+R_xS_{x'}+S_xR_{x'}\bigr].
\end{equation}
For example, the one-shell choice $x=x'=(1,t)$ has $S_x=1+(Q-1)t^2$ and
$R_x=2t+(Q-2)t^2$, and its gain over the hard window is
\begin{equation}\label{fw:one-shell-gain}
 \log\frac{(k+1)S_x^2+2S_xR_x}{k+1}-2(1+\delta)\log\bigl(1+(Q-1)t^p\bigr).
\end{equation}
Since $p>1$, its right derivative at $t=0$ is $4/(k+1)>0$, so a
sufficiently small first shell always improves on the hard window. The
certificate uses the exact finite functional \eqref{fw:product-shells}.

\subsection{Finite Periods and Lattice Counting}

Let $\mathcal P$ be a finite set of finite places of $F$ that split in
$K$. For $v\in\mathcal P$ write $v\mathcal O_K=\mathfrak P_v\,\iota\mathfrak
P_v$, where we also write $v$ for the prime ideal of $F$. The completion
of $K$ at $\mathfrak P_v$ is $F_v$; let $j_v\colon K\to F_v$ be the
completion map. The map $x\mapsto(j_v(x),j_v(\iota x))$ records the
completions at $\mathfrak P_v$ and at $\iota\mathfrak P_v$, and $\iota$
exchanges the two coordinates. If $N_{K/F}\beta=1$, then
$j_v(\beta)j_v(\iota\beta)=1$. Let $\mathcal O_v$, $\varpi_v$ and $Q_v$
be the valuation ring, a uniformizer and the residue cardinality of
$F_v$. Put
\[
 \mathbb A_{\mathcal P}=\prod_{v\in\mathcal P}(F_v\times F_v),
\]
with the Haar measure that gives $\prod_v(\mathcal O_v\times\mathcal
O_v)$ measure one, and let $\Gamma=\mathcal O_{K,\mathcal P}$ be the
ring of elements of $K$ that are integral at every prime other than the
$\mathfrak P_v$ and $\iota\mathfrak P_v$. We embed $\Gamma$ in
$\C^d\times\mathbb A_{\mathcal P}$ by
$\gamma\mapsto(\gamma_\infty,\gamma_f)$, where
$\gamma_f=(j_v(\gamma),j_v(\iota\gamma))_v$. A subgroup of $\mathbb
A_{\mathcal P}$ of the form
\[
 \mathcal C_f=\prod_{v\in\mathcal P}\bigl(\varpi_v^{a_v}\mathcal
 O_v\times\varpi_v^{b_v}\mathcal O_v\bigr),\qquad a_v,b_v\in\Z,
\]
is called \emph{rectangular}; its measure is
$|\mathcal C_f|=\prod_vQ_v^{-a_v-b_v}$.

\begin{lemma}\label{fw:lattice}
Let $\mathcal C_f$ be rectangular and put
$\mathfrak L=\prod_v\mathfrak P_v^{a_v}(\iota\mathfrak P_v)^{b_v}$.
\begin{enumerate}
\item[(a)] The elements $\gamma\in\Gamma$ with $\gamma_f\in\mathcal C_f$
form $\mathfrak L$, and $N\mathfrak L=|\mathcal C_f|^{-1}$, so
$\covol(\mathfrak L_\infty)=\covol_0/|\mathcal C_f|$.
\item[(b)] Every coset of $\mathcal C_f$ in $\mathbb A_{\mathcal P}$
contains $\gamma_f$ for some $\gamma\in\Gamma$.
\item[(c)] $\Gamma$ is a lattice in $\C^d\times\mathbb A_{\mathcal P}$.
If $P_{\mathfrak L}$ is a fundamental parallelepiped of
$\mathfrak L_\infty$, then $P_{\mathfrak L}\times\mathcal C_f$ is a
fundamental domain for $\Gamma$, of measure $\covol_0$.
\end{enumerate}
\end{lemma}

\begin{proof}
(a) The valuation of $j_v(\gamma)$ is the exponent of $\mathfrak P_v$
in $\gamma$, and that of $j_v(\iota\gamma)$ is the exponent of
$\iota\mathfrak P_v$. Together with integrality at the other primes,
$\gamma_f\in\mathcal C_f$ says exactly that $\gamma\in\mathfrak L$. As
$N\mathfrak P_v=N(\iota\mathfrak P_v)=Q_v$, we get
$N\mathfrak L=\prod_vQ_v^{a_v+b_v}=|\mathcal C_f|^{-1}$, and the
covolume follows from the formula of Section~\ref{geo:section}.

(b) Choose $a'_v\le a_v$ and $b'_v\le b_v$ such that the given coset
lies in $\mathcal C'_f=\prod_v(\varpi_v^{a'_v}\mathcal O_v\times
\varpi_v^{b'_v}\mathcal O_v)$, and let $\mathfrak L'\supset\mathfrak L$
be the corresponding ideal. The map $\gamma\mapsto\gamma_f$ induces a
homomorphism $\mathfrak L'/\mathfrak L\to\mathcal C'_f/\mathcal C_f$. It
is injective by (a), applied to $\mathcal C_f$. Both groups are finite
of the same order: $|\mathcal C'_f/\mathcal C_f|=|\mathcal
C'_f|/|\mathcal C_f|$, and
$|\mathfrak L'/\mathfrak L|=N\mathfrak L/N\mathfrak L'$
\cite[Proposition~4.2]{MilneANT}, which is the same number by (a). So
the map is bijective, and every coset of $\mathcal C_f$ in $\mathcal
C'_f$ contains $\gamma_f$ for some $\gamma\in\mathfrak L'\subset\Gamma$.
No principality is needed.

(c) Given $(z,c)\in\C^d\times\mathbb A_{\mathcal P}$, by (b) there is
$\gamma\in\Gamma$ with $c-\gamma_f\in\mathcal C_f$, and then a unique
$\gamma'\in\mathfrak L$ with $z-(\gamma+\gamma')_\infty\in P_{\mathfrak
L}$. If $(z,c)$ had two such representations, the difference of the
two elements of $\Gamma$ would have finite part in $\mathcal C_f$, hence
lie in $\mathfrak L$ by (a), and then it would vanish because
$P_{\mathfrak L}$ is a fundamental domain. Thus
$P_{\mathfrak L}\times\mathcal C_f$ is a fundamental domain, of measure
$(\covol_0/|\mathcal C_f|)\cdot|\mathcal C_f|=\covol_0$. Finally $\Gamma$
is discrete: for a bounded open set $Y\ni0$ in $\C^d$, its intersection
with the neighborhood $Y\times\mathcal C_f$ of $0$ consists of the
finitely many $\gamma\in\mathfrak L$ with $\gamma_\infty\in Y$.
\end{proof}

For $h\in\R^c$ let $\Gamma_h=\{(W_h\gamma_\infty,\gamma_f):\gamma\in
\Gamma\}$. It is again a lattice, with fundamental domains
$W_hP_{\mathfrak L}\times\mathcal C_f$ of measure $\covol_0$. For a
rectangular $\mathcal C_f$ put
\[
 H_f=d^{-1}\log|\mathcal C_f|.
\]
By Lemma~\ref{fw:lattice}(a), $N\mathfrak L=e^{-dH_f}$, and
Lemma~\ref{geo:dual} gives the dual product bound $\mu_f^d$ for every
$W_h\mathfrak L_\infty$, with $\mu_f=2e^{H_f/2-\ell}$. Accordingly, the
Fourier condition used in this section is
\begin{equation}\label{fw:fourier-condition}
 \sigma\cdot2e^{H_f/2-\ell}\ge\log M+2\log5+2 .
\end{equation}

\begin{lemma}\label{fw:counting}
Let $(f_\R,g_\C)$ be admissible with Fourier constants $(M,\sigma)$, and
let $\Psi_\infty=\prod_vf_\R^p\prod_jg_\C^p$ on $\C^d$. Let
$\Psi_f\ge0$ be a function on $\mathbb A_{\mathcal P}$ that is invariant
under translation by a rectangular group $\mathcal C_f$ and vanishes
outside a finite union of its cosets. If \eqref{fw:fourier-condition}
holds, with $H_f$ defined by this $\mathcal C_f$, then for every
$h\in\R^c$ and every $z_0\in\C^d\times\mathbb A_{\mathcal P}$,
\begin{equation}\label{fw:weighted-count}
 \sum_{\gamma\in\Gamma_h+z_0}\Psi_\infty(\gamma_\infty)\Psi_f(\gamma_f)
 \le\frac2{\covol_0}\int\Psi_\infty\int\Psi_f ,
\end{equation}
where $\gamma_\infty$ and $\gamma_f$ are the two components of $\gamma$.
\end{lemma}

\begin{proof}
Group the points by the coset of $\mathcal C_f$ containing their finite
component. If two points of $\Gamma_h+z_0$ have finite components in
the same coset, their difference comes from an element of $\mathfrak L$,
by Lemma~\ref{fw:lattice}(a). So the archimedean components of the
points in one coset form a translate of $W_h\mathfrak L_\infty$, or the
empty set. By Lemma~\ref{geo:poisson}, with the dual bound $\mu_f$ and
the Fourier bound for $\Psi_\infty$ from Section~\ref{geo:section}, the
sum of $\Psi_\infty$ over such a translate is at most
$2|\mathcal C_f|\int\Psi_\infty/\covol_0$. The values of $\Psi_f$ on its
finitely many nonzero cosets sum to $\int\Psi_f/|\mathcal C_f|$.
\end{proof}

If $\alpha\in\mathbb A_{\mathcal P}$ has coordinates
$(\varpi_v^{n_v},\varpi_v^{-n_v})$, multiplication by $\alpha$ preserves
the Haar measure and maps rectangular groups to rectangular groups of
the same measure. So \eqref{fw:fourier-condition} is unchanged, and
Lemma~\ref{fw:counting} applies after such a rescaling of $\Psi_f$,
even though the rescaling does not preserve $\Gamma$.

\subsection{Class Selection and Unit Averaging}

For $u\in\R^c$ and $n=(n_v)\in\Z^{\mathcal P}$ let
$\nu(u,n)\in\C^d\times\mathbb A_{\mathcal P}$ have archimedean part
$\nu(u)$ and finite part $(\varpi_v^{n_v},\varpi_v^{-n_v})_v$.

\begin{lemma}\label{fw:selection}
Let $\mathcal N\subset\Z^{\mathcal P}$ be finite, and let
$\mathcal W\colon\mathcal N\to[0,\infty)$ have positive sum. There are
$n_0\in\mathcal N$ with $\mathcal W(n_0)>0$, a subset
$\mathcal N'\subset\mathcal N$ with
\[
 \sum_{n\in\mathcal N'}\mathcal W(n)\ge\frac1{\kappa h_{\rm rel}}
 \sum_{n\in\mathcal N}\mathcal W(n),
\]
and elements $\beta_n\in\Gamma$, for $n\in\mathcal N'$, with
$N_{K/F}\beta_n=1$ and
\[
 \beta_n\mathcal O_K=\prod_v\mathfrak P_v^{\,n_v-n_{0,v}}
 (\iota\mathfrak P_v)^{-(n_v-n_{0,v})}.
\]
The cosets $\beta_n\mathcal O_K^1$ are pairwise disjoint, and
$\beta_n\mathcal O_K^1$ is the set of all elements of relative norm one
that generate this ideal.
\end{lemma}

\begin{proof}
Map $n$ to the class of $\mathfrak A(n)=\prod_v\mathfrak P_v^{n_v}$ in
$\Cl(K)$ modulo the image of $\Cl(F)$, a group of order $\kappa h_{\rm
rel}$. Let $\mathcal N'$ be a fiber of largest total weight, and choose
$n_0\in\mathcal N'$ with $\mathcal W(n_0)>0$. For $n\in\mathcal N'$ write
$\mathfrak A(n)\mathfrak A(n_0)^{-1}=x_n\mathfrak a_n\mathcal O_K$ with
$x_n\in K^\times$ and $\mathfrak a_n$ a fractional ideal of $F$, and put
$\beta_n=x_n/\iota x_n$. As in the proof of Lemma~\ref{geo:selection},
$N_{K/F}\beta_n=1$ and $\beta_n$ generates the displayed ideal, which
is supported on the primes $\mathfrak P_v,\iota\mathfrak P_v$; hence
$\beta_n\in\Gamma$. Two elements of relative norm one generating the
same ideal differ by an element of $\mathcal O_K^1$, and distinct labels
give distinct ideals.
\end{proof}

If $\beta\in\beta_n\mathcal O_K^1$, then $j_v(\beta)$ has valuation
$n_v-n_{0,v}$ and $j_v(\iota\beta)=j_v(\beta)^{-1}$. Let
$\varpi_0\in\mathbb A_{\mathcal P}$ have coordinates
$(\varpi_v^{n_{0,v}},\varpi_v^{-n_{0,v}})$. Then $\varpi_0\beta_f$ has
coordinates $(\varpi_v^{n_v}\omega_v,\varpi_v^{-n_v}\omega_v^{-1})$ with
units $\omega_v\in\mathcal O_v^\times$.

\subsection{The Finite-Place Transfer}

\begin{proposition}\label{fw:transfer}
Let $0<\delta<1$ and $C\in\R$. Let $(K_j/F_j)_{j\ge1}$ be quadratic
extensions satisfying (G1)--(G3) with a common $\lambda$, with
$d_j=[F_j:\Q]\to\infty$ and $\log L_{F_j}(1)\le d_jC$. For each $j$ let
$\mathcal P_j$ be a finite set of finite places of $F_j$ that split in
$K_j$, and for $v\in\mathcal P_j$ let $g_v$ be a shell profile on
$F_v^2$ of type $(Q_v,k_v,w^{(v)})$, where the types range over a fixed
finite set independent of $j$. Let $\mathcal F_{\delta,v}=\mathcal
F_\delta(g_v)$, let $(f_\R,g_\C)$ be admissible with Fourier constants
$(M,\sigma)$, and suppose that \eqref{fw:fourier-condition} holds for
every $j$, with $\mathcal C_f=\prod_{v\in\mathcal P_j}(\mathcal
O_v\times\varpi_v^{-k_v}\mathcal O_v)$, that is, with
$H_f=d_j^{-1}\sum_{v\in\mathcal P_j}k_v\log Q_v$. Put
\begin{equation}\label{fw:margin}
 \begin{split}
 \mathcal M_{\rm fin}(\theta)={}&
 \frac1d\sum_{v\in\mathcal P}\log\mathcal F_{\delta,v}
 -(1/2-\delta)\ell-C+(1-\delta)\log2+(1-\theta)\log\pi\\
 &+(1-2\theta)J_\R+\theta J_\C ,
 \end{split}
\end{equation}
evaluated for $K_j/F_j$. If these margins have a positive lower bound
over $j$, then there are finite sets $U_j\subset\R^2$ with
$|U_j|\to\infty$ and $u(U_j)/|U_j|^{1+\delta}\to\infty$.
\end{proposition}

\begin{proof}
Fix $\eta>0$ with $4\eta$ less than the lower bound of the margins, and
consider one extension $K/F$ of the sequence, with $d$ large.

\emph{The profile and its window.} On $\C^d\times\mathbb A_{\mathcal
P}$ let $f=\prod_vf_\R\prod_jg_\C\prod_{v\in\mathcal P}g_v$, with one
factor for each archimedean block and each finite place, let
$V=-\log f$ on the set where $f>0$, and put
\[
 A=\int f^p=A_\R^bA_\C^c\prod_{v\in\mathcal P}A(g_v),\qquad
 Z=I_\R^bI_\C^c\prod_{v\in\mathcal P}Z(g_v).
\]
At a finite place, the \emph{endpoint law} is the probability density
$g_v(z)g_v(z+s(n,\omega))/Z(g_v)$ on $F_v^2\times\Z\times\mathcal
O_v^\times$; the energy $-\log g_v$ of its first endpoint takes finitely
many values. The map $(z,n,\omega)\mapsto(-z-s(n,\omega),n,\omega)$
preserves this law and carries the first endpoint to minus the second,
and $g_v$ is even. Together with the archimedean endpoint laws of
Section~\ref{geo:section}, the product of all these laws is a
probability law under which the energies $V$ of the two endpoints have
the same distribution and are sums of independent terms. Let $d\bar m$
be their mean. The residue characteristics of the places in $\mathcal
P$ lie in a fixed finite set, so $|\mathcal P|=O(d)$, and since the
types range over a finite set, the variance of either energy is $O(d)$.
As in the proof of Proposition~\ref{geo:transfer}, with probability at
least $1/2$ for large $d$, both energies lie in
$[d(\bar m-\eta),d(\bar m+\eta)]$.

Let $\Omega$ be the set of points of $\C^d\times\mathbb A_{\mathcal P}$
where $f>0$ and $V\le d(\bar m+\eta)$. It is a bounded Borel set: the
archimedean energies are coercive and bounded below, and the finite
part lies in the compact support of $\prod_vg_v$, on which the finite
energies are bounded. It is invariant under rotations of the
archimedean coordinates, under $(x_v,y_v)\mapsto(\omega x_v,\omega'y_v)$
for units at each $v\in\mathcal P$, and under translation by
$\mathcal C_f$. For $u\in\R^c$ and $n\in\Z^{\mathcal P}$ let
$\Phi(u,n)=|\Omega\cap(\Omega-\nu(u,n))|$, and let
$\mathcal W(n)=\int_{\R^c}\Phi(u,n)\,du$. By the unit invariance, the
overlap of $\Omega$ with its translate by a vector whose archimedean
part has the moduli of $\nu(u)$ and whose finite part is
$(\varpi_v^{n_v}\omega_v,\varpi_v^{-n_v}\omega_v^{-1})_v$ equals
$\Phi(u,n)$. Only finitely many $n$ have $\mathcal W(n)\ne0$. Undoing the
endpoint density on the event of the previous paragraph, as in the
proof of Proposition~\ref{geo:transfer}, gives
\[
 \sum_{n}\mathcal W(n)\ge\tfrac12Ze^{2d(\bar m-\eta)},
 \qquad
 |\Omega|\le Ae^{pd(\bar m+\eta)} .
\]

\emph{Class selection and rescaling.} Apply Lemma~\ref{fw:selection} to
the finitely many $n$ with $\mathcal W(n)>0$, obtaining $n_0$,
$\mathcal N'$ and $\beta_n$, and let $\varpi_0$ be as above. Put
$\Omega'=\varpi_0^{-1}\Omega$, where $\varpi_0$ acts on the finite part.
Then $|\Omega'|=|\Omega|$, and $\Omega'$ is invariant under the
rectangular group $\mathcal C'_f=\varpi_0^{-1}\mathcal C_f$, of measure
$|\mathcal C_f|$. For $n\in\mathcal N'$, $\beta\in\beta_n\mathcal O_K^1$
and $h\in\R^c$, the element $\beta_h=(W_h\beta_\infty,\beta_f)\in
\Gamma_h$ satisfies
\[
 |\Omega'\cap(\Omega'-\beta_h)|=|\Omega\cap(\Omega-(W_h\beta_\infty,
 \varpi_0\beta_f))|=\Phi(l(\beta)-h,n),
\]
by the last paragraph of the previous subsection and the unit
invariance.

\emph{Points and edges.} For $h\in\R^c$ and $z_0$ in the fundamental
domain $W_hP_{\mathfrak L'}\times\mathcal C'_f$ of $\Gamma_h$ given by
Lemma~\ref{fw:lattice}, where $\mathfrak L'$ corresponds to
$\mathcal C'_f$, let $\mathcal X(h,z_0)=(\Gamma_h+z_0)\cap\Omega'$ and
$N(h,z_0)=|\mathcal X(h,z_0)|$, and let $\mathcal E(h,z_0)$ count the
pairs $(x,\beta)$ with $x\in\mathcal X(h,z_0)$,
$\beta\in\bigcup_{n\in\mathcal N'}\beta_n\mathcal O_K^1$ and
$x+\beta_h\in\Omega'$. Parametrize $z_0$ by $y\in[0,1)^{2d}$ through a
basis of $\mathfrak L'_\infty$, as in Section~\ref{geo:section}, and by
the Haar probability on $\mathcal C'_f$. Since
$\mathbf 1_{\Omega'}\le e^{pd(\bar m+\eta)}(f\circ\varpi_0)^p$ and
$(f\circ\varpi_0)^p$ is the product of an archimedean factor and a
function invariant under $\mathcal C'_f$, Lemma~\ref{fw:counting} gives
\[
 N(h,z_0)\le N_{\max}=\frac{2Ae^{pd(\bar m+\eta)}}{\covol_0}
\]
for all $(h,z_0)$. By tiling, the average of $N(h,z_0)$ over $z_0$ is
$|\Omega'|/\covol_0$, and the average of $\mathcal E(h,z_0)$ is
$\covol_0^{-1}\sum_{n\in\mathcal N'}\sum_{\beta\in\beta_n\mathcal O_K^1}
\Phi(l(\beta)-h,n)$. Averaging over $h\in\Pi^1$ as in
Lemma~\ref{geo:unfold-lemma}, the average of $\mathcal E$ becomes
\[
 \frac{w_K}{\operatorname{Reg}^1\covol_0}\sum_{n\in\mathcal N'}\mathcal W(n)
 \ge\frac{w_K}{\kappa h_{\rm rel}\operatorname{Reg}^1\covol_0}
 \sum_n\mathcal W(n).
\]
Measurability and integrability hold as in
Lemma~\ref{geo:joint-average-lemma}, since $\mathcal E\le N^2\le
N_{\max}^2$. The argument of that lemma yields a pair $(h,z_0)$ with
$N=N(h,z_0)>0$ and
\[
 \frac{\mathcal E(h,z_0)}N\ge\frac{w_K}{\kappa h_{\rm rel}
 \operatorname{Reg}^1}\,\frac{\sum_n\mathcal W(n)}{|\Omega|}
 \ge\frac{w_K}{\kappa h_{\rm rel}\operatorname{Reg}^1}\,
 \frac{Ze^{2d(\bar m-\eta)}}{2Ae^{pd(\bar m+\eta)}} .
\]
With $N\le N_{\max}$ and $p(1+\delta)=2$,
\[
 \frac{\mathcal E(h,z_0)}{N^{1+\delta}}\ge
 \frac{w_K\covol_0^\delta}{2^{1+\delta}\kappa h_{\rm rel}
 \operatorname{Reg}^1}\,\frac{Z}{A^{1+\delta}}\,e^{-4\eta d}.
\]
By \eqref{geo:mass-density} and $\log L_F(1)\le dC$, the first factor
is at least
$2^{-2-\delta}2^{d(1-\delta)}\pi^{(1-\theta)d}\lambda^{-(1/2-\delta)d}e^{-dC}$,
and $Z/A^{1+\delta}=\exp\bigl(bJ_\R+cJ_\C+\sum_v\log\mathcal
F_{\delta,v}\bigr)$. Hence
$\mathcal E/N^{1+\delta}\ge2^{-2-\delta}e^{d(\mathcal M_{\rm
fin}(\theta)-4\eta)}$.

\emph{The planar set.} Let $v_0$ be a real place of $F$ and let $U$ be
the set of $v_0$-coordinates of the points of $\mathcal X(h,z_0)$.
Distinct points differ by $(W_h\gamma_\infty,\gamma_f)$ with
$0\ne\gamma\in\Gamma$, whose $v_0$-coordinate is nonzero, so $|U|=N$.
Each counted pair $(x,\beta)$ gives two points at distance
$|\sigma_{v_0}(\beta)|=1$, and distinct pairs give distinct ordered
pairs of points. As in the proof of Proposition~\ref{geo:transfer},
$u(U)\ge\mathcal E/2$, the ratio $u(U)/|U|^{1+\delta}$ tends to
infinity, and the bound $u(U)\le|U|^2/2$ forces $|U|\to\infty$.
\end{proof}

\subsection{Uniform Local Types}
Suppose that $K/F$ has uniform local types $(e_r,f_r)_{r\in\mathcal R}$
in the sense of Section~\ref{geo:section}, that $\mathcal P$ is the set
of all places of $F$ above $\mathcal R$, and that at every place above
$r$ we use a shell profile of the same type $(r^{f_r},k_r,w^{(r)})$.
Write $\mathcal F_{\delta,r}$ for its functional. Since there are
$d/(e_rf_r)$ places above $r$,
\begin{equation}\label{fw:uniform}
 \frac1d\sum_{v\in\mathcal P}\log\mathcal F_{\delta,v}
 =\sum_{r\in\mathcal R}\frac{\log\mathcal F_{\delta,r}}{e_rf_r},
 \qquad
 H_f=\sum_{r\in\mathcal R}\frac{k_r\log r}{e_r}=H .
\end{equation}
In particular \eqref{fw:fourier-condition} coincides with
\eqref{geo:poisson-condition}, whatever the shell weights. For hard
windows, Lemma~\ref{fw:shell-lemma} gives
$\log\mathcal F_{\delta,r}=\log(k_r+1)-\delta k_rf_r\log r$, and the sum
in \eqref{fw:uniform} is $J-\delta H$. Thus
Proposition~\ref{fw:transfer} contains Proposition~\ref{geo:transfer},
and the shell weights replace the term $J-\delta H$ of
\eqref{geo:general-margin} by the finite functional without changing the
Fourier condition.
